\section{Global residue duality for the actual original-zeta
quotient}\label{global-residue-duality-for-the-actual-original-zeta-quotient}

24 September 2026. Independent derivation. Proof locators GZR0--GZR9.

\subsection{GZR0. Original objects, sources, and operation
domains}\label{gzr0.-original-objects-sources-and-operation-domains}

This calculation constructs a continuous, nondegenerate bilinear pairing
on the entire original quotient. It uses the original denominator
\(\zeta(s)\), every actual zero multiplicity, and the quotient Fréchet
topology. Its definition is a pair of absolutely convergent integrals.
An unrestricted infinite sum of residues is not used as its definition
or assumed to converge.

The controlling source notation remains \(Z_0,Z_1,Z_2,\tau\). Primitive
\(\tau\) is neither added nor assigned parity, an arithmetic coordinate,
or a numerical weight. All additions and derivatives below occur in the
existing complex receiving spaces of the already reconstructed original
zeta function. A receiving supported zero \(z_\lambda\) is not primitive
\(\tau\).

Before interpreting the result, the complete user passages
USR-9ad1c0a2d09dba92, USR-f55d16feb948d8c2, and USR-6152e3bc6302258c
were reread in the preserved corpus, together with
\nolinkurl{CORPUS_AND_OPERATION_RULES.md}. They concern the whole reconstructed
arithmetic system, admissible global transformations, and the requested
comparison with Deligne's weight argument. This note does not replace
that source by an independently invented spectrum.

The exact prior programme inputs are:

\begin{enumerate}
\def\labelenumi{\arabic{enumi}.}
\tightlist
\item
  \href{ORIGINAL_MELLIN_SPECTRAL_SYNTHESIS.md}{OMS1--OMS6, Original
  Mellin spectral synthesis}: the original test spaces, topological
  Mellin comparison, exact Schwartz image, uniform rapid division, full
  zero jets, and original primary sections.
\item
  RZ1--RZ8,
  \href{https://github.com/KokunoYumeto/zeta-function-research-reader/blob/main/workbenches/splitzero-tandem/continuations/20260924-cohomology-weight-and-character-lifts/gct-weight-control/proofs/ORIGINAL_ZETA_RESOLVENT_AND_PROJECTORS.md}{Original-zeta
  resolvent and projectors}: actual global representatives of each
  individual multiplicity jet and the full arithmetic action. These
  sections and the retained multiplier were read in
  \nolinkurl{quantum_tau_programme_bridge_20260924/sources/programme/cc_sheaf/ORIGINAL_ZETA_RESOLVENT_AND_PROJECTORS.md}.
  The linked retained copy is byte-identical: SHA256
  \nolinkurl{f3553e1362b4c2c1c1556098b55ee153be9fc716a9d4c35baedcf39bd28ad755}. The
  representatives are prior results, not a new construction claimed
  here.
\item
  The earlier
  \href{https://github.com/KokunoYumeto/zeta-function-research-reader/blob/4ba9285b66af58a6d58fc502a494c1fb45ca5b79/workbenches/tau-base-cohomology/TAU_BASE_MODEL.md\#L265-L322}{finite
  residue pairing and Jacobian trace contraction, TAU\_BASE\_MODEL §6}.
  The locally retained §6 was read. That calculation uses a finite
  completed-germ denominator. Its finite residue mechanism is prior
  work; GZR1--GZR8 establish the present global construction on the
  actual quotient with the original denominator. Only §6's receiving
  finite residue and trace algebra is used: that older document's
  introductory absorber interpretation of tau is not the current
  primitive identity source and is not imported. The current source
  remains the user's B1--B5 and retraction R1.
\end{enumerate}

The human analytic source underlying OMS is Ralf Meyer,
\href{https://arxiv.org/abs/math/0412277v3}{A spectral interpretation
for the zeros of the Riemann zeta function, arXiv:math/0412277v3},
original author TeX, function spaces and Zeta operator through the exact
Fourier--Laplace range theorem. The retained original author source is
sources/Meyer\_0412277v3/Meyer.tex, SHA256
\nolinkurl{ab9bc31f3c105a64fdc6f3b65ad16701dd8bf218fc90df98fe8d7f3c7c58c00e}. The
stronger rapid-division result invoked below is proved in OMS3; it is
not attributed without proof to Meyer's theorem.

Write \(\mathscr Z\) for the distinct actual nontrivial zeros of the
original \(\zeta\), and \(m_\rho\) for their actual positive
multiplicities. No assumption of simple zeros, RH, or positivity is
made.

\subsection{GZR1. The entire original quotient and the needed division
theorem}\label{gzr1.-the-entire-original-quotient-and-the-needed-division-theorem}

Retain \[
\begin{aligned}
S&=\{f\in\mathcal S(\mathbb R;\mathbb C):
 f(-v)=f(v),\ f(0)=0,\ \int_{\mathbb R}f(v)\,dv=0\},\\
\mathcal A&=\left\{k\in C^\infty(\mathbb R_{>0}):
 \sup_{u>0}(u^N+u^{-N})|(u\partial_u)^jk(u)|<\infty
 \text{ for every }N,j\ge0\right\},\\
\mathcal Ef(u)&=u^{1/2}\sum_{n\ge1}f(nu),\\
\mathcal Mk(s)&=\int_0^\infty k(u)u^{s-1/2}\frac{du}{u}.
\end{aligned}
\tag{GZR1.1}
\] The receiving Fréchet algebra is \[
\mathcal B=\{F\text{ entire}:\
 b_{A,M}(F):=\sup_{|\sigma|\le A,\ t\in\mathbb R}
 (1+|t|)^M|F(\sigma+it)|<\infty
 \text{ for all integers }A,M\ge0\}.
\tag{GZR1.2}
\] Products and multiplication by every fixed polynomial are continuous
on \(\mathcal B\), as follows immediately by distributing the displayed
weights. Reflection \(RF(s)=F(1-s)\) is continuous, since \[
b_{A,M}(RF)\le b_{A+1,M}(F).
\tag{GZR1.3}
\]

The exact original closed image theorem proved in OMS gives \[
\mathcal M(\mathcal ES)=\mathcal I,\qquad
\mathcal I=\{F\in\mathcal B:F^{(j)}(\rho)=0\
(\rho\in\mathscr Z,\ 0\le j<m_\rho)\},
\] \[
Q=\mathcal A/\mathcal ES
\ \xrightarrow[\mathcal M]{\ \cong\ }\ 
\mathcal Q=\mathcal B/\mathcal I .
\tag{GZR1.4}
\] The equality uses the actual image, not only its closure. The arrow
is a topological isomorphism. The ideal is closed because all the finite
derivative evaluations are continuous. We write \(\pi F=[F]\), and use
the quotient seminorms \[
\bar b_{A,M}([F])=\inf_{H\in\mathcal I}b_{A,M}(F+H).
\tag{GZR1.5}
\]

We use the following proved part of OMS3, with its original Schwartz
inverse. If \(H\in\mathcal I\), then \(H/\zeta\) extends holomorphically
throughout a neighborhood of the strip \(-1\le\Re s\le2\), and, for
every \(M\), its absolute value times \((1+|\Im s|)^M\) is bounded
uniformly on this strip by finitely many seminorms of \(H\). This is a
conclusion of the exact image theorem, not an assumption that
meromorphic division preserves \(\mathcal B\). Globally \(H/\zeta\) can
have the retained trivial-zero poles.

More explicitly, the actual inverse is \[
f_H(v)=\frac1{2\pi}\int_{\mathbb R}
\frac{H(2+it)}{\zeta(2+it)}v^{-2-it}\,dt,
\quad v>0,
\qquad f_H\in S,\quad \mathcal M\mathcal Ef_H=H .
\tag{GZR1.6}
\] Its continuation toward zero retains \[
\operatorname{Res}_{s=-2r}\frac{H(s)}{\zeta(s)}
 =\frac{H(-2r)}{\zeta'(-2r)},\qquad
\zeta'(-2r)=(-1)^r2^{-2r-1}\pi^{-2r}(2r)!\zeta(1+2r),
\tag{GZR1.7}
\] and its endpoint values retain \[
(H/\zeta)(0)=-2H(0),\qquad
(H/\zeta)(1)=0,\qquad
(H/\zeta)'(1)=H(1).
\tag{GZR1.8}
\] The uniform strip estimate is OMS3.8, proved by the full original
functional equation, the Hadamard product with all factors, and the
Gaussian-damped maximum principle before removal of that damping.
Equations (GZR1.6)--(GZR1.8) are the original Schwartz return used here.

The complete original functional multiplier is \[
\zeta(s)=\chi(s)\zeta(1-s),\qquad
\chi(s)=2^s\pi^{s-1}\sin(\pi s/2)\Gamma(1-s)
=\pi^{s-1/2}\frac{\Gamma((1-s)/2)}{\Gamma(s/2)}.
\tag{GZR1.9}
\] It is holomorphic and nonzero at every nontrivial zero. Thus
\(1-\rho\) is an actual zero of the same multiplicity, and
\(R\mathcal I=\mathcal I\).

\subsection{GZR2. Two absolutely convergent integrals and their
continuity}\label{gzr2.-two-absolutely-convergent-integrals-and-their-continuity}

For \(F,G\in\mathcal B\), define \[
\mathfrak B_\zeta(F,G)=\frac1{2\pi i}
\left(
\int_{2-i\infty}^{2+i\infty}
-\int_{-1-i\infty}^{-1+i\infty}
\right)\frac{F(s)G(1-s)}{\zeta(s)}\,ds .
\tag{GZR2.1}
\] Both displayed vertical integrals are oriented upward. Their
difference agrees with the vertical part of a counterclockwise
rectangle: the right side is upward and the left side is downward.

Here are absolute bounds using the original function, without a
replacement denominator. The reciprocal Euler series gives \[
|\zeta(2+it)^{-1}|\le\zeta(2).
\tag{GZR2.2}
\] For \(s=-1+it\), set \(y=t/2\). The full Gamma factors in (GZR1.9)
give \[
|\chi(-1+it)|^2
=\pi^{-3}y\coth(\pi y)(y^2+\tfrac14).
\tag{GZR2.3}
\] The continuous value at \(y=0\) is understood. To verify this
identity, use \[
\Gamma(-\tfrac12+iy)
=\frac{\Gamma(\tfrac12+iy)}{-\tfrac12+iy},
\quad
|\Gamma(1-iy)|^2=\frac{\pi y}{\sinh(\pi y)},
\quad
|\Gamma(\tfrac12+iy)|^2=\frac{\pi}{\cosh(\pi y)}
\] in the Gamma ratio in (GZR1.9). Since \[
y^2+\tfrac14\ge\frac{(1+|t|)^2}{8},
\qquad
y\coth(\pi y)\ge\frac{1+|t|}{2\pi},
\] the same reciprocal Euler series at \(1-s=2-it\) yields \[
|\zeta(-1+it)^{-1}|
\le4\pi^2\zeta(2)(1+|t|)^{-3/2}.
\tag{GZR2.4}
\] For the second elementary inequality, put \(x=\pi|y|\). Then
\(x\coth x\ge\max(1,x)\ge(1+2x/\pi)/2\); the limiting value at zero is
one.

All four numerator evaluations on the two edges lie in \(|\Re s|\le2\).
Since \(\int_{\mathbb R}(1+|t|)^{-2}\,dt=2\), equations (GZR2.2) and
(GZR2.4) prove the explicit bound \[
|\mathfrak B_\zeta(F,G)|
\le C_\zeta\, b_{2,1}(F)b_{2,1}(G),
\qquad
C_\zeta=\frac{\zeta(2)(1+4\pi^2)}{\pi}.
\tag{GZR2.5}
\] Thus each integral is absolutely convergent and the pairing is
jointly continuous and complex bilinear on the whole of \(\mathcal B\).

There is no endpoint pole in the integrand: \(\zeta(0)=-1/2\), while
\(1/\zeta(s)\) has a simple zero at \(s=1\). The trivial zeros
\(-2,-4,\ldots\) lie outside this strip; they remain in the original
source return (GZR1.7). No contour enclosing them has been substituted
for (GZR2.1).

\subsection{GZR3. Descent in both slots through the actual Schwartz
image}\label{gzr3.-descent-in-both-slots-through-the-actual-schwartz-image}

Suppose first \(F\in\mathcal I\). The function \(F/\zeta\) is
holomorphic and uniformly rapidly decreasing on the closed strip by
GZR1. The product \[
K(s)=(F(s)/\zeta(s))G(1-s)
\] has the same properties. On the rectangle with vertical sides at
\(-1,2\) and horizontal sides at imaginary parts \(\pm T\), Cauchy's
theorem gives zero for its complete contour integral. The horizontal
sides have length three and tend uniformly to zero as \(T\to\infty\).
The absolutely convergent vertical sides therefore give \[
\mathfrak B_\zeta(F,G)=0\qquad(F\in\mathcal I).
\tag{GZR3.1}
\] This proof invokes the actual division theorem and its Schwartz
return, not formal cancellation of the zeta denominator.

If \(G\in\mathcal I\), then \(RG\in\mathcal I\), with all
multiplicities, by (GZR1.9). Apply the identical argument to \[
K(s)=F(s)\bigl(G(1-s)/\zeta(s)\bigr).
\] It gives \[
\mathfrak B_\zeta(F,G)=0\qquad(G\in\mathcal I).
\tag{GZR3.2}
\] There is consequently a well-defined bilinear form on the entire
original quotient: \[
B_\zeta:\mathcal Q\times\mathcal Q\longrightarrow\mathbb C,
\qquad B_\zeta([F],[G])=\mathfrak B_\zeta(F,G).
\tag{GZR3.3}
\] For every \(H,K\in\mathcal I\), descent and (GZR2.5) give the same
bound with \(F+H,G+K\). Taking the two independent infima proves \[
|B_\zeta(x,y)|
\le C_\zeta\,\bar b_{2,1}(x)\bar b_{2,1}(y).
\tag{GZR3.4}
\] In particular the descended bilinear form is jointly continuous for
the actual quotient topology.

\subsection{GZR4. Every local multiplicity coefficient and its finite
perfect
form}\label{gzr4.-every-local-multiplicity-coefficient-and-its-finite-perfect-form}

Fix an actual zero \(\rho\), put \(m=m_\rho\), and write \(t=s-\rho\).
Retain the full local expansion \[
\zeta(\rho+t)=t^m u_\rho(t),\qquad
u_\rho(t)=\sum_{k\ge0}u_{\rho,k}t^k,\qquad
u_{\rho,k}=\frac{\zeta^{(m+k)}(\rho)}{(m+k)!},
\quad u_{\rho,0}\ne0.
\tag{GZR4.1}
\] Define the exact reciprocal coefficients \[
c_{\rho,0}=u_{\rho,0}^{-1},\qquad
c_{\rho,n}=-u_{\rho,0}^{-1}
\sum_{r=1}^{n}u_{\rho,r}c_{\rho,n-r}\quad(n\ge1).
\tag{GZR4.2}
\] Equating coefficients in \(u_\rho(t)\sum_{n\ge0}c_{\rho,n}t^n=1\)
proves that these are precisely the Taylor coefficients of
\(u_\rho^{-1}\). They retain every original zeta derivative contributing
to the corresponding order.

The exact residue is \[
\operatorname{Res}_{s=\rho}
\frac{F(s)G(1-s)}{\zeta(s)}\,ds
=
\sum_{\substack{i,j,k\ge0\\i+j+k=m-1}}
\frac{F^{(i)}(\rho)}{i!}
\frac{G^{(j)}(1-\rho)}{j!}
(-1)^j c_{\rho,k}.
\tag{GZR4.3}
\] Indeed \(G(1-\rho-t)=\sum_{j\ge0}G^{(j)}(1-\rho)(-t)^j/j!\), and the
residue is the coefficient of \(t^{m-1}\) in the product of this series,
the Taylor series of \(F\), and \(u_\rho^{-1}\). This proves the
displayed formula, including its sign and every factorial.

The local algebras are the actual full primary algebras \[
A_\rho=\mathbb C[t]/(t^m),\qquad
A_{1-\rho}=\mathbb C[v]/(v^m),\quad v=s-(1-\rho).
\tag{GZR4.4}
\] For \(P=\sum_{i<m}p_it^i\) and \(Q=\sum_{j<m}q_jv^j\), their local
pairing is \[
\beta_\rho(P,Q)
=[t^{m-1}]\left(
P(t)Q(-t)\sum_{k=0}^{m-1}c_{\rho,k}t^k
\right).
\tag{GZR4.5}
\] Its matrix in the ordered monomial bases is \[
M_{\rho,ij}=
\begin{cases}
(-1)^j c_{\rho,m-1-i-j},&i+j\le m-1,\\
0,&i+j>m-1,
\end{cases}
\qquad 0\le i,j<m.
\tag{GZR4.6}
\] Reversing the columns yields a triangular matrix. The column-reversal
sign is \((-1)^{m(m-1)/2}\), and the product of the anti-diagonal signs
is also \((-1)^{m(m-1)/2}\). Hence \[
\det M_\rho=c_{\rho,0}^{\,m}
=\left(\frac{m!}{\zeta^{(m)}(\rho)}\right)^m\ne0.
\tag{GZR4.7}
\] Every multiplicity jet is therefore retained. For example, when
\(m=2\), the actual matrix is \[
\begin{pmatrix}c_{\rho,1}&-c_{\rho,0}\\
c_{\rho,0}&0\end{pmatrix},
\quad
c_{\rho,0}=\frac{2}{\zeta''(\rho)},\quad
c_{\rho,1}=-\frac{2\zeta'''(\rho)}{3\zeta''(\rho)^2}.
\tag{GZR4.8}
\] These are exact finite calculations, not an assumption that the
global pairing is a convergent infinite sum of them.

\subsection{GZR5. Actual global isolators and a justified finite-residue
theorem}\label{gzr5.-actual-global-isolators-and-a-justified-finite-residue-theorem}

We recall the existing RZ representatives with all factors. Set \[
f_*(v)=\frac\pi2v^2(2\pi v^2-3)e^{-\pi v^2},\qquad
F_*(s)=\mathcal M\mathcal Ef_*(s)
=\frac{s(s-1)}8\pi^{-s/2}\Gamma(s/2)\zeta(s).
\tag{GZR5.1}
\] This is an actual auxiliary transform in \(\mathcal I\); it does not
replace the denominator in (GZR2.1). Its full exceptional values are \[
F_*(0)=F_*(1)=\frac18,\qquad
F_*(-2r)=F_*(1+2r)
=\frac{(1+2r)2r}{8}\pi^{-(1+2r)/2}
\Gamma((1+2r)/2)\zeta(1+2r)\ne0.
\tag{GZR5.2}
\] For \(\eta\in\mathscr Z\), put \(m=m_\eta\), \[
U_\eta(s)=\frac{F_*(s)}{(s-\eta)^m},\quad
a_{\eta,k}=\frac{F_*^{(m+k)}(\eta)}{(m+k)!},
\] \[
d_{\eta,0}=a_{\eta,0}^{-1},\quad
d_{\eta,n}=-a_{\eta,0}^{-1}
\sum_{r=1}^{n}a_{\eta,r}d_{\eta,n-r},\quad
e_\eta(s)=U_\eta(s)\sum_{n=0}^{m-1}d_{\eta,n}(s-\eta)^n .
\tag{GZR5.3}
\] Removal of the zero at \(\eta\) proves \(U_\eta\) entire. Division by
the fixed polynomial outside a bounded neighborhood, and Taylor division
inside, prove \(U_\eta\in\mathcal B\). The finite polynomial factors
therefore give \(e_\eta\in\mathcal B\).

The reciprocal recursion proves \[
e_\eta(s)=1+O((s-\eta)^m),\qquad
E_{\eta,j}(s)=e_\eta(s)(s-\eta)^j,\quad 0\le j<m.
\] \[
\frac{E_{\eta,j}^{(r)}(\nu)}{r!}
=
\begin{cases}1,&\nu=\eta,\ r=j,\\0,&\text{otherwise},\end{cases}
\quad \nu\in\mathscr Z,\quad 0\le r<m_\nu.
\tag{GZR5.4}
\] At every other zero the full original vanishing order is retained by
\(F_*\). Thus these equations hold at all zeros simultaneously. Each
function has its actual original representative \[
k_{\eta,j}(u)=\frac1{2\pi}\int_{\mathbb R}
E_{\eta,j}(\tfrac12+it)u^{-it}\,dt\in\mathcal A.
\tag{GZR5.5}
\] The Mellin isomorphism in GZR1 proves the asserted membership and
identity of its transform.

For completeness, all derivatives used in the reciprocal recursion
retain the original multiplier \[
C(s)=\frac18s(s-1)\pi^{-s/2}\Gamma(s/2),
\] \[
a_{\eta,k}=\sum_{h=0}^k
\frac{C^{(h)}(\eta)}{h!}
\frac{\zeta^{(m+k-h)}(\eta)}{(m+k-h)!},
\] \[
C^{(h)}(s)=\frac{\pi^{-s/2}}8
\sum_{\substack{a+b+c=h\\0\le a\le2}}
\frac{h!}{a!b!c!}(s(s-1))^{(a)}
\left(-\frac{\log\pi}{2}\right)^b2^{-c}\Gamma^{(c)}(s/2).
\tag{GZR5.6}
\]

Here is the contour result needed to apply these actual representatives.
Suppose \(F,G\in\mathcal B\), and the meromorphic function
\(F(s)G(1-s)/\zeta(s)\) has nonremovable poles inside the strip at only
a finite set \(S_0\subset\mathscr Z\). Put \[
P_{S_0}(s)=\prod_{\rho\in S_0}(s-\rho)^{m_\rho}.
\tag{GZR5.7}
\] The product \(H(s)=P_{S_0}(s)F(s)G(1-s)\) lies in \(\mathcal B\). At
every zero outside \(S_0\), pole removability says that the numerator
already has the required full zero order. At a zero in \(S_0\), the
explicit factor supplies it. Thus \(H\in\mathcal I\).

GZR1 now proves \(H/\zeta\) uniformly rapidly decreasing on the whole
strip. At sufficiently large \(|\Im s|\), the polynomial \(P_{S_0}\) is
bounded below by a positive constant times
\((1+|\Im s|)^{\deg P_{S_0}}\), uniformly in that strip. Consequently \[
\frac{F(s)G(1-s)}{\zeta(s)}
=\frac{H(s)/\zeta(s)}{P_{S_0}(s)}
\] also tends uniformly to zero faster than any inverse power on the
horizontal sides of sufficiently tall rectangles. The residue theorem
and absolute convergence of the vertical edges give \[
B_\zeta([F],[G])
=\sum_{\rho\in S_0}
\operatorname{Res}_{s=\rho}
\frac{F(s)G(1-s)}{\zeta(s)}\,ds .
\tag{GZR5.8}
\] This is a proved finite sum. It does not use estimates of \(1/\zeta\)
close to infinitely many unrelated zeros.

In particular, if \(G=E_{1-\rho,j}\), its reflected numerator vanishes
to every full zero order except possibly at \(\rho\); (GZR5.8) has only
that one residue. The same is true with the first slot supported by an
\(E_{\rho,j}\). This is the exact bridge from the global contour pairing
to the perfect local form.

\subsection{GZR6. Global nondegeneracy, exact radicals, and primary
adjoints}\label{gzr6.-global-nondegeneracy-exact-radicals-and-primary-adjoints}

Take a nonzero \(x=[F]\in\mathcal Q\). Synthesis (GZR1.4) implies that
at least one actual full jet is nonzero. Fix such a \(\rho\), and let
\(\ell<m=m_\rho\) be the least index with \[
f_\ell=F^{(\ell)}(\rho)/\ell!\ne0.
\] Choose \(G=E_{1-\rho,m-1-\ell}\). Equations (GZR5.8) and (GZR4.3)
give exactly \[
B_\zeta(x,[G])
=(-1)^{m-1-\ell} f_\ell c_{\rho,0}\ne0.
\tag{GZR6.1}
\] Indeed the only possible term has \(i=\ell,j=m-1-\ell,k=0\): all
smaller \(i\) vanish, while larger \(i\) exceed the permitted total
degree.

For a nonzero \(y=[G]\), choose a zero \(\eta\) with first nonzero
Taylor coefficient \(g_\ell=G^{(\ell)}(\eta)/\ell!\), and put
\(\rho=1-\eta\). Taking \(F=E_{\rho,m_\eta-1-\ell}\) gives \[
B_\zeta([F],y)=(-1)^\ell g_\ell c_{\rho,0}\ne0.
\tag{GZR6.2}
\] Thus both radicals of the global form vanish. Equivalently, the two
radicals of the original form \(\mathfrak B_\zeta\) on \(\mathcal B\)
are exactly \(\mathcal I\), not larger subspaces.

Let \(s_\rho:A_\rho\to\mathcal Q\) be the existing primary section
supplied by (GZR5.4), and \(j_\rho\) the full Taylor-jet map. The actual
projector is \(P_\rho=s_\rho j_\rho\). Formula (GZR5.8) proves \[
B_\zeta(s_\rho P,s_\eta Q)=0\quad(\eta\ne1-\rho),
\] \[
B_\zeta(s_\rho P,s_{1-\rho}Q)=\beta_\rho(P,Q),
\qquad
B_\zeta(P_\rho x,y)=B_\zeta(x,P_{1-\rho}y).
\tag{GZR6.3}
\] For the last identity, both sides are the same local residue at
\(\rho\); the projectors retain exactly the jets appearing in (GZR4.3).
These equations prove orthogonality and the reflected adjoint relation
for actual primary subspaces of the original quotient, with all
nilpotents retained.

\subsection{GZR7. Full real action and the reflected companion
denominator}\label{gzr7.-full-real-action-and-the-reflected-companion-denominator}

For every \(a>0\), retain \[
W_ak(u)=a^{1/2}k(u/a),\qquad
T_aF(s)=a^sF(s),\qquad
\mathcal MW_a=T_a\mathcal M.
\tag{GZR7.1}
\] Multiplication by \(a^s\) is continuous on every strip and preserves
all zero orders, so it acts on the entire quotient. Inside each
integral, \[
(T_aF)(s)(T_aG)(1-s)
=a^sa^{1-s}F(s)G(1-s)=aF(s)G(1-s).
\] Therefore the full global similitude is \[
B_\zeta(T_ax,T_ay)=a B_\zeta(x,y).
\tag{GZR7.2}
\] No factor \(a^{1/2}\) has been absorbed into a different spectral
parameter.

Let \(L[F]=[sF]\), corresponding on the original source to
\(1/2-u\partial_u\). Directly in the absolutely convergent integrals, \[
B_\zeta(Lx,y)=B_\zeta(x,(1-L)y),\qquad
B_\zeta(Lx,y)+B_\zeta(x,Ly)=B_\zeta(x,y).
\tag{GZR7.3}
\] More generally \(B_\zeta(P(L)x,y)=B_\zeta(x,P(1-L)y)\) for every
complex polynomial \(P\). At paired primary blocks this retains the full
nilpotent reflection \(t\mapsto-v\), not only the scalar eigenvalues.

Define the companion using its actual original denominator: \[
B_{\zeta^\vee}([F],[G])=
\frac1{2\pi i}
\left(
\int_{2-i\infty}^{2+i\infty}
-\int_{-1-i\infty}^{-1+i\infty}
\right)
\frac{F(s)G(1-s)}{\zeta(1-s)}\,ds .
\tag{GZR7.4}
\] Its edge bounds exchange (GZR2.2) and (GZR2.4), so both integrals
converge absolutely. Substitute \(w=1-s\) in each integral before
subtraction. A single upward integral on \(\Re s=c\) becomes an upward
integral on \(\Re w=1-c\): the reversed traversal and \(ds=-dw\) cancel.
The two vertical edges then exchange places, yielding the remaining
minus sign \[
\boxed{\ B_\zeta(y,x)=-B_{\zeta^\vee}(x,y).\ }
\tag{GZR7.5}
\] This identity proves descent and continuity of the companion as well.
Since \(R^2=1\) and \((RF)(s)(RG)(1-s)=G(s)F(1-s)\), it also proves \[
B_\zeta(Rx,Ry)=B_\zeta(y,x)=-B_{\zeta^\vee}(x,y).
\tag{GZR7.6}
\] The sign is an orientation calculation, not a choice of a phase.

The full functional equation records the exact denominator comparison \[
\frac1{\zeta(1-s)}
=\frac{2^s\pi^{s-1}\sin(\pi s/2)\Gamma(1-s)}{\zeta(s)}
=\frac{\pi^{s-1/2}\Gamma((1-s)/2)}
{\Gamma(s/2)\zeta(s)}.
\tag{GZR7.7}
\] These are meromorphic equalities with their removable values
retained. At \(s=0\) the companion reciprocal has a zero; at \(s=1\) it
equals \(1/\zeta(0)=-2\). The Gamma and endpoint behavior is not
discarded. Equations (GZR7.5)--(GZR7.7) do not assert that \(B_\zeta\)
itself is symmetric, skew-symmetric, Hermitian, or positive.

\subsection{GZR8. Continuous injections into the strong dual and their
weak-*
density}\label{gzr8.-continuous-injections-into-the-strong-dual-and-their-weak--density}

Let \(\mathcal Q'\) denote the continuous complex-linear dual, and let
\(\mathcal Q'_b\) have the strong topology of uniform convergence on
bounded subsets of the actual quotient Fréchet space. Define \[
\mathcal D_L(x)(y)=B_\zeta(x,y),\qquad
\mathcal D_R(y)(x)=B_\zeta(x,y).
\tag{GZR8.1}
\] Equation (GZR3.4) proves that both are well-defined maps into the
continuous dual. GZR6 proves both injective.

For a bounded subset \(K\subset\mathcal Q\), its associated strong
seminorm is \[
p_K(\lambda)=\sup_{y\in K}|\lambda(y)|.
\] Since every continuous seminorm is bounded on a bounded set, (GZR3.4)
gives \[
p_K(\mathcal D_Lx)
\le C_\zeta\left(\sup_{y\in K}\bar b_{2,1}(y)\right)
\bar b_{2,1}(x).
\tag{GZR8.2}
\] The coefficient is finite. This proves continuity into the strong
dual for every defining strong seminorm. The same calculation with the
slots exchanged proves continuity of \(\mathcal D_R\).

Both images are dense for the weak-* topology
\(\sigma(\mathcal Q',\mathcal Q)\). Here is a direct proof retaining the
full function space. Fix \(\lambda\in\mathcal Q'\) and finitely many
\(y_1,\ldots,y_n\in\mathcal Q\). Consider the linear map \[
A:\mathcal Q\to\mathbb C^n,\qquad
A(x)=(B_\zeta(x,y_1),\ldots,B_\zeta(x,y_n)).
\tag{GZR8.3}
\] If a vector of coefficients \((c_1,\ldots,c_n)\) annihilates its
range, then \[
B_\zeta\left(x,\sum_{j=1}^n c_jy_j\right)=0
\quad\text{for all }x.
\] Right nondegeneracy implies \(\sum_jc_jy_j=0\), and hence
\(\sum_jc_j\lambda(y_j)=0\). In finite-dimensional linear algebra the
range of \(A\) is the common kernel of all linear functionals
annihilating it: extend a basis of its range to a basis of
\(\mathbb C^n\) to verify this statement. Thus \[
(\lambda(y_1),\ldots,\lambda(y_n))\in A(\mathcal Q).
\tag{GZR8.4}
\] One element \(\mathcal D_Lx\) therefore agrees exactly with
\(\lambda\) on every prescribed finite list. This proves weak-* density.
The proof for \(\mathcal D_R\) uses left nondegeneracy.

The action on this dual is the exact character-\(a\)-twisted
contragredient action \[
(T_a^\diamond\lambda)(y)=a\,\lambda(T_{a^{-1}}y).
\tag{GZR8.5}
\] A character with value \(a\) is not itself a proof that
\(\mathcal Q\) has Deligne weight one. When \(a=q\) is a Frobenius norm,
that one-dimensional character has Deligne weight two, since its
absolute value is \(q^{2/2}\). No purity conclusion is assigned to the
receiving space by this notation. A continuous linear automorphism
\(T_{a^{-1}}\) sends bounded sets to bounded sets, so this is continuous
on the strong dual. Formula (GZR7.2) gives \[
\mathcal D_LT_a=T_a^\diamond\mathcal D_L,\qquad
\mathcal D_RT_a=T_a^\diamond\mathcal D_R.
\tag{GZR8.6}
\] For the first equality, \(B_\zeta(T_ax,y)=aB_\zeta(x,T_{a^{-1}}y)\);
the second follows by exchanging the slots.

These are continuous injective maps with weak-* dense images. Neither
surjectivity onto the continuous dual nor a topological self-duality has
been proved or assumed.

\subsection{GZR9. All receiving support labels and the exact scope of
the
result}\label{gzr9.-all-receiving-support-labels-and-the-exact-scope-of-the-result}

Let \(L_{\mathrm{supp}}\) be the existing support lattice with top
element \(1_{L_{\mathrm{supp}}}\). For a receiving complex vector space
\(V\), retain \[
G_{L_{\mathrm{supp}}}(V)
=\{(0,\lambda):\lambda\in L_{\mathrm{supp}}\}
\ \cup\
\{(v,1_{L_{\mathrm{supp}}}):v\in V\}.
\tag{GZR9.1}
\] For every linear map \(f:V\to W\), its actual support-preserving lift
is \[
G_{L_{\mathrm{supp}}}(f)(v,\lambda)=(f(v),\lambda).
\tag{GZR9.2}
\] It is well-defined because a nonzero amplitude occurs only at top
support. Identity and composition are preserved by direct substitution.
In particular the quotient, original Mellin comparison, full primary
projectors, reflection, real action, and the two dual injections all
have these lifts. A zero output retains its original \(\lambda\). If
\(f\) is injective, then equality of its lifted outputs first gives
equality of labels and then equality of amplitudes, so the lift is
injective.

The bilinear form itself admits a lift retaining both input labels: \[
\widetilde B_\zeta:
G_{L_{\mathrm{supp}}}(\mathcal Q)\times
G_{L_{\mathrm{supp}}}(\mathcal Q)
\longrightarrow G_{L_{\mathrm{supp}}\times L_{\mathrm{supp}}}(\mathbb C),
\] \[
\widetilde B_\zeta((x,\lambda),(y,\mu))
=(B_\zeta(x,y),(\lambda,\mu)).
\tag{GZR9.3}
\] If either input has nontop support, its amplitude is zero, so the
output amplitude is zero as required. If the output amplitude is
nonzero, both inputs have top support, giving the top pair in the
product lattice. This proves well-definedness without merging either
input label. All scalar and action identities above preserve these pairs
of labels. This is an explicit map on the stated receiving carrier; an
ordinary tensor product that identifies all zero amplitudes is not
substituted for it.

The global calculation has therefore produced a full-multiplicity
continuous duality map for the actual
\(\mathcal B/\mathcal I=\mathcal A/\mathcal ES\), together with its
exact real action, original denominator, reflected transpose, and
strong-dual continuity. The finite residue forms are its proved
restrictions, not a replacement for the whole space. The result gives no
positive-definite form and no estimate forcing \(\Re\rho=1/2\). It
supplies the nondegenerate global pairing on which an additional proved
weight or positivity argument would have to act; it is not itself that
argument.

\subsection{Explicit global continuation: finite lifting and a nonzero
remainder}\label{explicit-global-continuation-finite-lifting-and-a-nonzero-remainder-occurrence-3}

The earlier uncertainty about algebraic surjectivity is now resolved by
\nolinkurl{GLOBAL_POLYNOMIAL_DUAL_DEFECT.md} PGD1--PGD6, independently checked in
PGC0--PGC7: the actual algebraic cokernel is nonzero. Its explicit class
c\_zeta is the original-zeta contour functional with the constant
arithmetic function 1 in the first slot. A lift would require a rapidly
decreasing entire representative to equal 1 at all actual zeros, which
the proved unbounded zero heights exclude. PGD constructs an injected
rational-function line, entire representatives retaining all Hermite
pole corrections, and the whole real/prime orbit. This is a receiving
cokernel class, not primitive tau or a claimed off-critical zero.

SCL0--SCL9 and DPL0--DPL9 prove every nonzero polynomial in the actual
cokernel generator invertible. Thus all finite-dimensional invariant
dual spaces lift uniquely, continuously and with the original real/prime
action. DPL gives explicit extension splittings and computes every
degree of the cyclic derived Hom into the actual cone. These results
retain its H and H′ terms, all faithful closed copies and support
labels. The weak-* Hausdorff quotient remains zero alongside the now
proved nonzero algebraic quotient. Current source Z\_0, Z\_1/tau and
Z\_2 are unchanged. Complete proofs and exact human citations are in the
linked continuation files; the full geometric weight-separation goal
remains active.
